\section{System}
\label{sec:system}

We study an open-source Google Sheets add-on in Apps Script (about 9{,}400 lines of JavaScript, 163 unit tests). A user types a request; the assistant plans, asks for approval if the plan writes, executes through a registry of typed tools, and writes into the workbook (Figure~\ref{fig:arch}).

\begin{figure*}[t]
\centering
\resizebox{0.97\textwidth}{!}{%
\begin{tikzpicture}[
  box/.style={draw, rounded corners=1.5pt, minimum height=9mm, minimum width=24mm, align=center, font=\small, inner sep=3pt},
  hl/.style={box, fill=blue!10, draw=blue!60!black},
  arr/.style={-{Latex[length=2mm]}, thick},
  lbl/.style={font=\footnotesize, fill=white, inner sep=1pt}
]
\node[box] (req) at (0,0) {Request};
\node[box] (ctx) at (3.1,0) {Active-sheet\\context};
\node[hl] (ret) at (6.4,0) {Budgeted\\retrieval};
\node[box] (plan) at (9.7,0) {Planner\\(approval if writes)};
\node[box] (agent) at (13.2,0) {Formula agent\\(LLM call)};
\node[hl] (ver) at (13.2,-2.2) {\gc\\parse + six layers};
\node[box] (tool) at (9.7,-2.2) {Tool registry\\(four gates)};
\node[box] (wb) at (6.4,-2.2) {Workbook\\+ undo stack};
\node[hl] (ssrf) at (9.7,-4.2) {URL validator\\+ redirect re-check};
\draw[arr] (req) -- (ctx);
\draw[arr] (ctx) -- (ret);
\draw[arr] (ret) -- (plan);
\draw[arr] (plan) -- (agent);
\draw[arr] (agent) -- (ver) node[lbl, midway, right=1mm] {candidate formula};
\draw[arr] (ver) -- (tool) node[lbl, midway, above] {valid};
\draw[arr] (tool) -- (wb) node[lbl, midway, above] {write};
\draw[arr] (ssrf) -- (tool) node[lbl, midway, right=1mm] {fetch only if allowed};
\draw[arr, dashed] (ver.east) -- ++(1.5,0) |- (agent.east) node[lbl, pos=0.25, right] {errors, warnings, hints (repair)};
\end{tikzpicture}}
\caption{Request flow. Shaded boxes are evaluated in this paper: retrieved cross-sheet context, the grounded verifier with its repair feedback, and the URL validator in front of the fetch tool.}
\label{fig:arch}
\end{figure*}


\subsection{Grounded verification (\gc)}
\label{sec:gc}
Each candidate formula is parsed into a typed AST by a small recursive-descent parser (the earlier verifier matched regular expressions against raw text) and checked in six deterministic layers that read only the open workbook:
\begin{enumerate}\itemsep0pt\topsep0pt
\item \textbf{Structural.} Leading \code{=}, balanced quotes and parentheses outside string literals, injection patterns, successful parse.
\item \textbf{Symbols.} Referenced sheets exist (case-insensitively, as in Sheets); a bare identifier is a function, named range or \code{LET}/\code{LAMBDA} variable (a header used as a name is rejected with the column's letters as a hint); function names within edit distance 1--2 of a real one are errors with ``did you mean'', other unknown names are warnings (custom functions); \code{QUERY} column letters exist.
\item \textbf{Bounds.} Each reference is checked against the extent of \emph{the sheet it points at}. A range wholly beyond the used columns is an error; one that merely extends past the last data row is headroom (a note); \code{COUNTA}, \code{COUNTBLANK} and \code{ISBLANK} may point at empty cells.
\item \textbf{Shape.} Argument counts for about 90 functions, lookup indices beyond the lookup range, literal \code{INDEX} positions, unequal range sizes in \code{SUMIFS}, \code{COUNTIFS}, \code{SUMIF}, \code{SUMPRODUCT}.
\item \textbf{Grounding.} Numeric aggregates over a text column; criteria and lookup keys that never occur in the column they search. The workbook proves neither impossible, so they are \emph{suspicious}: warnings with the column's real values as the hint, optionally sent to repair.
\item \textbf{Circularity.} The target cell lies inside a referenced range (\code{=SUM(F2:F81)} typed into F81, \code{=SUM(F:F)} into column F) or on a dependency chain through other formulas, where range dependencies are expanded to the formula cells inside them; \code{ROW}, \code{ROWS}, \code{COLUMN}, \code{COLUMNS} read only geometry and add no dependency.
\end{enumerate}
\begin{figure*}[!t]
\centering
\resizebox{\textwidth}{!}{%
\begin{tikzpicture}[
  box/.style={draw, rounded corners=1.5pt, align=left, font=\footnotesize, inner sep=3pt},
  okb/.style={box, draw=green!50!black, fill=green!10},
  badb/.style={box, draw=red!70!black, fill=red!10},
  wrnb/.style={box, draw=orange!80!black, fill=orange!15},
  hl/.style={box, fill=blue!8, draw=blue!60!black},
  arr/.style={-{Latex[length=1.8mm]}, thick},
  lbl/.style={font=\scriptsize, fill=white, inner sep=1pt}
]
% candidate and its syntax tree
\node[box, text width=5.7cm] (req) at (2.85,0) {\textbf{Request:} ``total revenue for product Flange''\\[1pt]\textbf{Candidate:} \texttt{=SUMIF(C2:C80,"Flangee",Revenue)}};
\node[box, minimum width=1.6cm, align=center] (f) at (2.85,-1.75) {\texttt{SUMIF}};
\node[okb, align=center, minimum width=1.6cm] (a1) at (0.7,-3.15) {\texttt{C2:C80}\\[-1pt]\scriptsize range};
\node[wrnb, align=center, minimum width=1.6cm] (a2) at (2.85,-3.15) {\texttt{"Flangee"}\\[-1pt]\scriptsize string};
\node[badb, align=center, minimum width=1.6cm] (a3) at (5.0,-3.15) {\texttt{Revenue}\\[-1pt]\scriptsize bare name};
\draw[arr] (req) -- (f);
\draw[thick] (f) -- (a1); \draw[thick] (f) -- (a2); \draw[thick] (f) -- (a3);
% layers
\node[box, draw=none, font=\footnotesize\bfseries] at (9.2,0.75) {Six layers read the open workbook};
\node[okb, text width=4.7cm] (l1) at (9.2,0) {\textbf{1~Structural} \hfill parses, balanced};
\node[badb, text width=4.7cm] (l2) at (9.2,-0.85) {\textbf{2~Symbols} \hfill \texttt{Revenue} is no name};
\node[okb, text width=4.7cm] (l3) at (9.2,-1.7) {\textbf{3~Bounds} \hfill within rows 2--80};
\node[okb, text width=4.7cm] (l4) at (9.2,-2.55) {\textbf{4~Shape} \hfill 3 arguments, equal sizes};
\node[wrnb, text width=4.7cm] (l5) at (9.2,-3.4) {\textbf{5~Grounding} \hfill criterion not in column};
\node[okb, text width=4.7cm] (l6) at (9.2,-4.25) {\textbf{6~Circularity} \hfill no cycle};
\draw[thick] (f.east) -- (6.3,-1.75);
\draw[thick] (6.3,0) -- (6.3,-4.25);
\foreach \y/\n in {0/l1,-0.85/l2,-1.7/l3,-2.55/l4,-3.4/l5,-4.25/l6} \draw[arr] (6.3,\y) -- (\n.west);
% messages and retry
\node[box, text width=6.1cm] (msg) at (15.7,-1.15) {\textbf{Error:} ``Revenue'' is not a function, named range or cell reference.\\[1pt]\textbf{Hint:} ``Revenue'' is a column header (column~F). Reference it by cell range, e.g.\ F2:F80, not by name.\\[1pt]\textbf{Suspicious:} criterion ``Flangee'' never occurs in C2:C80 (``Product''). Values in that column include: Sprocket, Widget, Gizmo, Doohickey, Flange, Gadget. Closest match: ``Flange''.};
\node[okb, text width=6.1cm] (fix) at (15.7,-4.25) {\textbf{Retry} (temperature 0.1) with the diagnostics appended:\\\texttt{=SUMIF(C2:C80,"Flange",F2:F80)}\\ all layers pass; the formula may be written.};
\draw[arr] (l2.east) -- ++(0.6,0) |- (msg.west);
\draw[arr] (l5.east) -- ++(0.6,0) |- (msg.west);
\draw[arr] (msg.south) -- (fix.north) node[lbl, midway, right] {invalid: retry};
\end{tikzpicture}}
\caption{Worked example. The candidate formula mistypes a criterion and uses a column header as a name. GroundCheck parses it, reads the workbook's sheets, headers and column values, rejects it, and returns the column letter and the real values as repair feedback; the second attempt passes. The messages are the verifier's own output for a generated workbook (\texttt{eval/worked\_example.js}).}
\label{fig:anatomy}
\end{figure*}

Errors make a formula \emph{invalid} and trigger a retry (Figure~\ref{fig:anatomy} walks through an example). The verifier cannot know whether a formula answers the question: \code{=SUM(Quantity)} over the wrong numeric column is valid and wrong. After a failed attempt the formula and the layers' errors, warnings and hints (``Column ``Revenue'' is column D'', ``values in that column include East, West, North'') are appended to the conversation; temperature decays from 0.2 to 0.1 to 0 over at most two retries.

\subsection{What a workbook can decide}
\label{sec:levels}
A candidate formula can be wrong in ways that differ in how much of the workbook is needed to see it (Table~\ref{tab:levels}). Whether it parses, whether the names it uses exist, whether its ranges lie inside the data, whether its arguments fit the function, and whether it sits on a cycle can all be decided from the workbook alone. Whether a criterion matches the data can be decided only as a \emph{suspicion}, since a value missing from a column may be intended. Whether a valid formula answers the user's question cannot be decided without the user's intent: \code{=SUM(Quantity)} over the wrong numeric column parses, names real things, fits, is acyclic and matches the data. \gc{} works on the first five levels and leaves the last to the user. The experiments probe exactly this boundary, and we use the distinction between \emph{loud} faults (the cell shows an error) and \emph{silent} ones (a plausible but wrong value) to measure it.

\begin{table}
\caption{What a workbook can decide about a candidate formula, from syntax to intent, with the share of \sheetfault{} faults of that kind that each verifier blocks/flags (\%). The first four levels are decidable from the workbook alone; the data level is decidable only as a suspicion; whether a valid formula answers the question is not decidable without the user's intent.}
\label{tab:levels}

\centering\footnotesize
\begin{adjustbox}{max width=\columnwidth}\begin{tabular}{>{\raggedright\arraybackslash}p{3.7cm} c r cc}
\toprule
\textbf{Level of correctness} & \textbf{Decidable} & \textbf{n} & \textbf{Shipped} & \textbf{\gc} \\
\midrule
Syntax: does it parse? & yes & 3{,}173 & 100/100 & \textbf{100/100} \\
Names and extent: do the sheets, functions, headers and ranges it uses exist, within the data? & yes & 3{,}423 & 29/81 & \textbf{90/99} \\
Shape: do the arguments fit the function? & yes & 468 & 0/27 & \textbf{100/100} \\
Dependencies: is the target cell on a cycle? & yes & 4{,}320 & 40/60 & \textbf{100/100} \\
Data: do criteria and types match the column? & in part & 948 & 0/25 & \textbf{0/100} \\
Intent: does it compute what was asked? & no & 1{,}024 & 0/32 & \textbf{4/4} \\
\bottomrule
\end{tabular}\end{adjustbox}
\end{table}


\subsection{Retrieval and the fetch boundary}
\label{sec:retrieval}
A deterministic analyzer splits the workbook into tables with inferred types and statistics, builds a formula dependency graph, and emits one chunk per table listing each column's \emph{letter}, type and statistics and the data-row span (a first version listed headers only; Section~\ref{sec:e5} shows why that failed). Seven signals score chunks against the request (active sheet, referenced by the active cell's formula, lexical overlap, shared headers, neighbouring table, graph proximity, conversation overlap) and a greedy knapsack fills a token budget. In the original add-on this context reached the planner but not the formula agent; we pass it on.

As a secondary component, the same idea of checking before acting applies where an LLM-chosen string leaves the add-on, the fetch tool, the setting of server-side request forgery (SSRF)~\cite{owaspssrf}. The original filter matched regular expressions against the whole URL. Our validator parses it and classifies the \emph{host}: IPv4 in any \code{inet\_aton} form, IPv6 including IPv4-mapped, NAT64 and 6to4, the reserved ranges of RFC~6890~\cite{rfc6890}, single-label and internal-only names, wildcard-DNS services, embedded credentials, and characters on which URL parsers disagree. Redirects are not auto-followed: every hop is re-validated. A static check cannot see DNS rebinding.
